\documentclass[oneside,a4paper,14pt]{extarticle}
\usepackage[a4paper,top=20mm,bottom=20mm,left=30mm,right=10mm]{geometry}
\usepackage{amsmath, amssymb, amsfonts, amsthm}
\usepackage[T2A]{fontenc}
\usepackage[utf8]{inputenc}
\usepackage[russian]{babel}
\usepackage{textcomp}
\usepackage{indentfirst}
\usepackage{graphicx}
\usepackage{float}
\usepackage{wrapfig}
\usepackage{caption}
\usepackage{tocloft}
\usepackage{seqsplit}
\allowbreak
\renewcommand{\cftsecleader}{\cftdotfill{\cftdotsep}}
\usepackage[breaklinks]{hyperref}
\usepackage{titlesec}
\usepackage{fancyvrb}
\usepackage{multirow}
\usepackage{xcolor}
\titleformat{\section}{\normalsize\bfseries}{\thesection}{1em}{}
\titleformat{\subsection}{\normalsize\bfseries}{\thesubsection}{1em}{}
\renewcommand\baselinestretch{1.33}\normalsize
\setlength{\parindent}{1.25cm}

\captionsetup[figure]{labelsep=period}

\begin{document}
\newpage\thispagestyle{empty}
\newpage
\setcounter{page}{1}
\tableofcontents
\newpage

\section*{Цель}
\addcontentsline{toc}{section}{Цель}
Разработать и исследовать управляющее устройство.

\section*{Задание}
\addcontentsline{toc}{section}{Задание}
Выполнить сдвиг 8-разрядного двоичного кода $X$ на $N$ разрядов в сторону старших разрядов (сдвиг логический).

На рис. \ref{fig:mem} показано распределение памяти микропрограмм по адресам памяти.
\begin{figure}[H]
	\centering
	\includegraphics[width=0.95\textwidth]{Рис1.png}
	\caption{Распределение памяти микропрограмм}
	\label{fig:mem}
\end{figure}

\newpage

\section*{Кодирование микропрограммы}
\addcontentsline{toc}{section}{Кодирование микропрограммы}
В табл. 1 приведено кодирование микропрограммы. Микроинструкция CONT ($\text{I3--I0} = 1110$) выполняется безусловно и передаёт управление следующей микрокоманде. Микроинструкция CJP ($\text{I3--I0} = 0011$) проверяет логическое условие (признак $Z$, выбираемый битами $\text{A1 A0} = 00$) и осуществляет переход по адресу, заданному на шине данных $\text{D11--D0}$. Микроинструкция JZ ($\text{I3--I0} = 0000$) выполняет безусловный переход к началу микропрограммы (адрес 0000).

Результат кодирования микропрограммы (при тестировании выхода из цикла с условием $Z=1$ в МК03) представлен в табл. 1.

\begin{flushright}
	Таблица 1
\end{flushright}
\begin{center}
	Кодирование микропрограммы
\end{center}
\begin{table}[H]
	\centering
	\renewcommand{\arraystretch}{1.3}
	\resizebox{\textwidth}{!}{
		\begin{tabular}{|c|c|c|c|c|c|c|c|c|c|c|c|c|}
			\hline
			№   & Шина данных & Код МИ  & \multicolumn{3}{c|}{Упр. признаками} & \multicolumn{3}{c|}{Упр. сигналы} & \multicolumn{4}{c|}{Признаки}                                                                       \\ \hline
			МК  & D11 - D0    & I3 - I0 & A1 A0                                & U                                 & \textasciicircum CCE          & C0 & \textasciicircum RLD & \textasciicircum OE & Z & F3 & OVR & C4 \\ \hline
			00: & 0000        & 1110    & 00                                   & 0                                 & 0                             & 1  & 1                    & 0                   & 0 & 0  & 0   & 0  \\ \hline
			01: & 0004        & 0011    & 00                                   & 1                                 & 0                             & 1  & 1                    & 0                   & 0 & 0  & 0   & 0  \\ \hline
			02: & 0000        & 1110    & 00                                   & 0                                 & 0                             & 1  & 1                    & 0                   & 0 & 0  & 0   & 0  \\ \hline
			03: & 0002        & 0011    & 00                                   & 0                                 & 0                             & 1  & 1                    & 0                   & 1 & 0  & 0   & 0  \\ \hline
			04: & 0000        & 0000    & 00                                   & 0                                 & 0                             & 1  & 1                    & 0                   & 0 & 0  & 0   & 0  \\ \hline
		\end{tabular}
	}
	\label{tab:microcommands}
\end{table}

\newpage

\section*{Пример}
\addcontentsline{toc}{section}{Пример}
В качестве первого примера рассмотрим случай $N = 0$. По алгоритму работы операционного устройства в микрокоманде МК00 (код МИ $1110$ — CONT) в регистр $R0$ записывается число $X$, после чего управление безусловно передаётся МК01. В микрокоманде МК01 в регистр $R9$ записывается количество сдвигов $N = 0$, при этом формируется признак нуля $Z = 1$.

В МК01 используется микроинструкция условного перехода CJP (код МИ $0011$), что соответствует проверке условия $Z = 1$. Поскольку количество сдвигов равно нулю ($Z = 1$), выполнение циклического участка не требуется. Микроинструкция CJP осуществляет переход из МК01 непосредственно к завершающей микрокоманде МК04 по адресу $0004$, заданному на шине данных $\text{D11--D0}$. Микрокоманды МК02 и МК03 пропускаются.

На рис. \ref{fig:before_mk01} показано состояние устройства управления перед выполнением микрокоманды МК01. Проверяется признак $Z = 1$.
\begin{figure}[H]
	\centering
	\includegraphics[width=0.75\textwidth]{До01.png}
	\caption{Состояние УУ перед выполнением МК01}
	\label{fig:before_mk01}
\end{figure}

На рис. \ref{fig:after_mk01} показано состояние устройства управления после выполнения МК01 при $Z = 1$. На выходной шине $Y$ формируется адрес $0004$, соответствующий микрокоманде МК04.
\begin{figure}[H]
	\centering
	\includegraphics[width=0.75\textwidth]{Из01В04.png}
	\caption{Переход из МК01 к МК04 при $Z=1$}
	\label{fig:after_mk01}
\end{figure}

В качестве второго примера рассмотрим случай $N = 1$ (соответствует настройке признаков в табл. 1 и на скриншоте микропрограммы: в МК01 задано $Z = 0$, а в МК03 задано $Z = 1$). После выполнения МК01 в регистр $R9$ записывается значение $1$. Поскольку $R9 \neq 0$, устанавливается признак $Z = 0$.

Условие перехода ($Z = 1$) в МК01 не выполняется, поэтому управление последовательно передаётся следующей микрокоманде МК02.

В МК02 (код МИ $1110$ — CONT) выполняется логический сдвиг числа $X$ в сторону старших разрядов на одну позицию, и управление передаётся МК03. В МК03 количество оставшихся сдвигов в регистре $R9$ уменьшается на единицу, а также выполняется микроинструкция CJP (код МИ $0011$) (проверка условия $Z = 0$ для возврата на адрес $0002$).

После уменьшения счётчика устанавливается признак $Z = 1$. Поскольку все необходимые сдвиги выполнены, условие возврата ($Z = 0$) к МК02 не выполняется. Управление последовательно передаётся МК04 (код МИ $0000$ — JZ), которая осуществляет безусловный переход к началу микропрограммы (МК00).

На рис. \ref{fig:before_mk03} показано состояние устройства управления перед выполнением МК03. После уменьшения счётчика проверяется признак $Z$.
\begin{figure}[H]
	\centering
	\includegraphics[width=0.75\textwidth]{До03.png}
	\caption{Состояние УУ перед выполнением МК03}
	\label{fig:before_mk03}
\end{figure}

На рис. \ref{fig:after_mk03} показано состояние устройства управления после выполнения МК03 при $Z = 1$. Условие возврата к МК02 не выполняется, поэтому управление передаётся МК04.
\begin{figure}[H]
	\centering
	\includegraphics[width=0.75\textwidth]{Из03В04.png}
	\caption{Переход из МК03 к МК04 при $Z=1$}
	\label{fig:after_mk03}
\end{figure}
\newpage
\section*{Вывод}
\addcontentsline{toc}{section}{Вывод}
В ходе лабораторной работы было разработано устройство управления с программируемой логикой для выполнения логического сдвига восьмиразрядного двоичного числа $X$ на $N$ разрядов в сторону старших разрядов. Была составлена микропрограмма, состоящая из пяти микрокоманд. Для управления последовательностью их выполнения использовались микроинструкции CONT, CJP и JZ.

Были реализованы условные переходы по признаку $Z$, позволяющие организовать циклическое выполнение сдвига и завершение операции при обнулении счётчика оставшихся сдвигов.

В результате экспериментальной проверки работы устройства управления было подтверждено корректное выполнение переходов между микрокомандами в зависимости от значения признака $Z$.

\end{document}